% Known-source correction only: replace ex:ch14-line-cycles and the
% following thm:current-fundamental-cycle with these complete blocks.
% The preceding ex:ch14-skeleton-current is not changed by this packet.

\begin{example}[Cycle currents on the Berkovich line]
\label{ex:ch14-line-cycles}
Return to $X=\mathbf P^{1,\An}_k$ with $k$ trivially valued and
algebraically closed, as in \Cref{sec:ch14-line-examples}.  Every analytic
$0$-cycle is a locally finite sum $Z=\sum_bm_b[b]$ over rigid points.
Since $X$ is compact, this sum is in fact finite.  Its integration current
is the measure
\[
 \delta_Z=\sum_bm_b\,\delta_{x_b},\qquad
 \langle\delta_Z,f\rangle=\sum_bm_bf(x_b),
 \qquad f\in A_{X,c}^{0,0}.
\]
Thus $\langle\delta_Z,1\rangle=\sum_bm_b$ is a legitimate pairing.

Every analytic $1$-cycle is $m[X]$, and its current is $m\,\delta_X$.
The current $\delta_X$ has bidimension $(1,1)$, hence bidegree $(0,0)$,
and acts on $(1,1)$-forms:
\[
                 \langle\delta_X,\alpha\rangle=\int_X\alpha,
                         \qquad \alpha\in A_{X,c}^{1,1}.
\]
It does not act on the constant function $1$.  In particular, its
integration cannot be identified with a \emph{mass against $1$} or with the
length of a chosen one-point model skeleton.

The calibrated rays in the star description of
\Cref{sec:ch14-line-examples} give a direct nonzero test.  Fix a branch
$b\in k$ and use $t=\log|T-b|\in(-\infty,0)$ on its interior.  Choose
$\chi\in C_c^\infty((-\infty,0))$ with $\chi\ge0$ and $\chi\not\equiv0$.
The form
\[
                        \alpha=\chi(t)\,\dprime t\wedge\ddprime t,
\]
extended by zero off that branch interval, is a smooth compactly
supported $(1,1)$-form on $X$.  The branch calibration has weight $1$, so
\[
                 \langle\delta_X,\alpha\rangle
                     =\int_{-\infty}^{0}\chi(t)\,\rd t>0.
\]
This is the actual star geometry over a trivially valued field;
see \cite[Example~6.7.3]{CLD25}.  A coarser choice of model skeleton
does not remove these calibrated ray tests.  Closedness of $\delta_X$
is \Cref{cor:ch14-delta-x-closed}.

Finally, \Cref{ex:ch14-line-model-function}, with $r=1$ and $c_1=1$,
gives the identity of currents
\[
                   \dprimeddprime[u]=\delta_{p_1}-\delta_\eta,
\]
where $u=\log\max\{|T-b_1|,\rme^{-1}\}$ on the branch $b_1$ and
$u=0$ elsewhere.  Here $p_1$ and $\eta$ are nonrigid points:
their Dirac masses are currents of bidimension $(0,0)$, but they are
not integration currents of analytic $0$-cycles.  This truncated-potential
identity must therefore be distinguished from the untruncated
Poincar\'e--Lelong formula \eqref{eq:appJ-cld-pl} for the analytic divisor of a meromorphic
function, treated in \Cref{chap:cld-poincare-lelong}.
\end{example}

\begin{theorem}
\label{thm:current-fundamental-cycle}
For every pure-dimensional closed analytic subspace $Y\subset X$,
\[
                         \delta_Y=\delta_{[Y]}.
\]
Thus nilpotent and embedded structure enters precisely through the generic
multiplicities of the fundamental cycle.
\end{theorem}
\begin{proof}
We give the calibrated-degree argument of
\cite[Proposition~16.6.3]{CLD25}.  Put $r=\dim Y$, let
$i:Y\hookrightarrow X$, and fix $\alpha\in A_{X,c}^{r,r}$.
Then $\omega=i^*\alpha$ is a smooth compactly supported top-degree
form on $Y$, and
\[
                         \langle\delta_Y,\alpha\rangle=\int_Y\omega.
\]

Let $(Y_j)_j$ be the locally finite family of reduced irreducible
components of $Y$, and let $m_j$ be the generic multiplicity of $Y$
along $Y_j$, so that $[Y]=\sum_jm_j[Y_j]$.
Let $Y^\circ$ be the Zariski-open subset consisting of points lying
on exactly one irreducible component.  The complement is the union
of the pairwise intersections of distinct components and has analytic
dimension less than $r$.  By
\cite[Proposition~15.2.1(b)]{CLD25}, the support of $\omega$ is
disjoint from that complement.  This is ambient vanishing near the
lower-dimensional analytic set, not merely vanishing of a restriction
to it.  In particular,
\[
                         \int_Y\omega=\int_{Y^\circ}\omega.
\]
We do not remove any component on which $Y$ is generically nonreduced.

Write $V_j=Y^\circ\cap |Y_j|$, with its induced open analytic subspace
structure from $Y$; this structure may still be nonreduced.
The $V_j$ are disjoint open subspaces of $Y^\circ$.
Only finitely many meet $\Supp\omega$, and $\omega_j=\omega|_{V_j}$
has compact support.  Let
\[
                       \nu_j:(V_j)_{\mathrm{red}}\longrightarrow V_j
\]
be the finite reduction immersion.  Its calibrated degree on the
generic locus relevant to integration is $1/m_j$, by
\cite[Proposition~13.3.7]{CLD25}.  Concretely, for an Artinian generic
local ring $A$ of length $m_j$, the calibrated-degree quotient is
\[
                  \frac{\operatorname{length}_A(A_{\mathrm{red}})}
                       {\operatorname{length}_A(A)}
                       =\frac{1}{m_j}.
\]
This is a calibrated degree; no flatness of the reduction immersion
or preservation of length under graded reduction is asserted.

The finite-map integration formula
\cite[Proposition~15.4.3]{CLD25}, applied to $\nu_j$ and $\omega_j$,
therefore gives
\[
             \int_{(V_j)_{\mathrm{red}}}\nu_j^*\omega_j
                          =\frac1{m_j}\int_{V_j}\omega_j.
\]
The same support argument identifies the left-hand integral with
$\int_{Y_j}i_j^*\alpha$, where $i_j:Y_j\hookrightarrow X$.
Consequently
\[
 \langle\delta_Y,\alpha\rangle
   =\sum_j\int_{V_j}\omega_j
   =\sum_jm_j\int_{Y_j}i_j^*\alpha
   =\langle\delta_{[Y]},\alpha\rangle.
\]
The sum is finite for this test form.  Since $\alpha$ was arbitrary,
the required equality follows.
\end{proof}
